% Section 14: industrial case studies.
\section{Industrial case studies}\label{sec:cases}

Benchmarks show that a solver is correct; case studies show that it fits the problems the problem
statement is about. \file{nirnay/cases/} builds \nCases{} optimisation models from public data, most of
it Indian. Every module carries a \code{CASE\_INFO} record of its sources (URL and licence), of what is
taken from them unchanged, and of what is assumed. \file{docs/CASE\_STUDIES.md} documents each model in
full. Each case is written to \file{data/cases/<case>.mps} and re-solved by HiGHS as a reference.
\Cref{tab:cases} gives the sizes.

\begin{table}[ht]
\centering\small
\caption[Case-study models]{Case-study models: problem class and size of the default instance, read
from the MPS files in \file{data/cases/} (integer columns in parentheses).}
\label{tab:cases}
\begin{tabularx}{\linewidth}{@{}X l l r r r@{}}
\toprule
Case & Module & Class & Rows & Columns & Nonzeros\\
\midrule
\rows{data/cases_sizes.tex}
\bottomrule
\end{tabularx}
\end{table}

\subsection{The models}
\paragraph{Refinery crude selection and blending (LP)}
The model chooses how much of each crude to run on each of MRPL's three crude distillation units
(capacities from MRPL's public project filing), where each distillation cut goes, and how to blend
BS-VI petrol (MS\,91, MS\,95), BS-VI diesel, ATF, LPG, naphtha and fuel oil to the BIS
specifications (IS~2796, IS~1460, IS~1571), to maximise gross margin at PPAC's 2024--25 average
prices. Crude assays come from the US Department of Energy's Strategic Petroleum Reserve (public
domain). Every cut of every crude is its own stream, so the property blending is linear. For a product
pool $P$ with components $i$ of mass $m_i$ and density $\rho_i$, a volumetric specification such as the
octane number is written
\begin{equation}\label{eq:blend}
  \sum_{i\in P}\frac{m_i}{\rho_i}\bigl(\mathrm{RON}_i-\mathrm{RON}_{\min}\bigr)\ \ge\ 0 ,
\end{equation}
which is linear in the mass flows. Density, cetane index, smoke point, naphthalenes and sulphur are
handled the same way.

\paragraph{Twelve-month production plan (MILP)}
Twelve monthly copies of the refinery network are linked by product inventories and storage limits.
Crude is bought in integer parcels, at most three crudes are processed per month (tank segregation),
and the FCC severity is a binary choice with at most two changeovers per year. Monthly demand follows
PPAC's national consumption pattern.

\paragraph{Thermal unit commitment (MILP)}
The 73 thermal units of the RTS-GMLC system over 48 hours, with the tight and compact formulation of
Morales-España, Latorre and Ramos~\cite{moralesespana2013} as published with PGLib-UC: on, start and
stop binaries, start-up categories, ramping, minimum up and down times, spinning reserve, and a
piecewise-linear production cost.

\paragraph{Economic dispatch with a DC network (QP)}
Quadratic generation cost over the 2000-bus PGLib-OPF case, with DC power flow, line limits and
angle-difference limits. The HiGHS optimum matches the DC objective PGLib-OPF publishes for the case.

\paragraph{Capacitated facility location (MILP)}
OR-Library's \code{capb} instance~\cite{beasley1990} (100 depots, 1000 customers) with the strong
linking constraints $x_{ij}\le y_i$. The HiGHS optimum equals the published optimum. It stands in for
depot and terminal network design, for which no Indian dataset with costs and capacities is public.

\paragraph{BS-VI fuel distribution across India (LP)}
A transport model from 21 Indian refineries to 36 states and union territories. It meets every state's
2024--25 MS and HSD sales (PPAC) with the least tonne-kilometres, using refinery and state locations
from Wikidata.

\paragraph{Crude-oil unloading and blending schedule (MILP)}
The crude scheduling problems of Lee et al.~\cite{lee1996} in the multi-operation sequencing
formulation of Mouret, Grossmann and Pestiaux~\cite{mouret2009} (the MILP relaxation that is the first
step of their method). All four instances reproduce the published optimal values.

\subsection{Reference solutions and \nirnay{}'s results}
\Cref{tab:cases-ref} gives the HiGHS reference for each case file, and \cref{tab:cases-runs} the
record of \nirnay{}'s runs kept in the case-study document.

\begin{table}[ht]
\centering\small
\caption[HiGHS reference results for the cases]{HiGHS \verHighs{} reference results for the case files
(\file{data/cases/reference.csv}). Refinery objectives are negative gross margins in thousand US\$.
The unit-commitment reference was solved to a \ucGapTarget\% gap target (final gap \ucGap\%), as in
PGLib-UC's own reference script, so it is not a proven optimum.}
\label{tab:cases-ref}
\begin{tabular}{@{}l l l r r r@{}}
\toprule
Case & Class & Status & Objective & MIP gap & Time (s)\\
\midrule
\rows{data/cases_ref.tex}
\bottomrule
\end{tabular}
\end{table}

\ifnum\haveCaseRuns=1
\begin{table}[ht]
\centering\small
\caption[\nirnay{} on the case studies]{\nirnay{} on the case studies (from the ``NIRNAY on the
cases'' record in \file{docs/CASE\_STUDIES.md}; one process per run, 120\,s limit, times including
Numba compilation on first call).}
\label{tab:cases-runs}
\begin{tabular}{@{}l l l r r r@{}}
\toprule
Case & Method & Status & Objective & Rel.\ diff. & Time (s)\\
\midrule
\rows{data/cases_runs.tex}
\bottomrule
\end{tabular}
\end{table}

\nCaseRunsOptimal{} of the \nCaseRuns{} recorded runs reach the HiGHS optimum: both LP cases with the
simplex and the interior point, and the economic dispatch QP. The MILP cases do not yet solve.
The twelve-month refinery plan stops at the time limit with a poor incumbent, the crude schedule
returns nothing within the wall limit, and unit commitment and facility location, the two largest
models in \cref{tab:cases}, were not attempted with branch-and-bound. This is the same gap as on MIPLIB
(\cref{sec:mip-results}), and closing it is the first item of the roadmap.
\fi
